\section{Exact-MoE-SubSpec}
\label{sec:method}

\subsection{One state hierarchy, two execution views}

\system targets batch-one inference under a ceiling $B$ that cannot hold the routed
BF16 expert pool. For layer $\ell$, $C^4_{\ell}$ and $C^{16}_{\ell}$ denote resident
W4 bases and materialized BF16 experts; both share the same calibrated device budget.
As Fig.~\ref{fig:overview}(b) shows, one round has four steps: the W4 view drafts up
to $\gamma$ tokens, logical rollback restores the clean prefix, the native BF16 view
verifies the root and proposals with on-demand exact experts, and only
target-authorized state is committed. The two views share the router, attention and
dense modules, prefill, and one physical KV allocation. Thus the system changes how
expert state is represented and supplied, not which target logits authorize output.

\subsection{One W4 expert, two roles}

Offline preparation quantizes every target expert's gate/up/down projections into a
group-64 W4 substitute. The W4 view preserves the target architecture, router,
routing rule, and all non-routed modules, yielding an aligned training-free draft
\cite{wang2025subspec}. Crucially, the same W4 object is also the base from which
verification reconstructs the exact target expert. The design therefore removes an
auxiliary routed draft store while turning draft residency into reusable verifier
state.

Quantization alone is approximate. For a dequantized-base BF16 pattern $b$ and target
pattern $w$, let $h(x)=x[15{:}7]$ and $m(x)=x[6{:}0]$. A correction stores $m(w)$ and
a two-bit relation code
\begin{equation}
c=\begin{cases}
0,&h(w)=h(b),\\
1,&h(w)=h(b)-1,\\
2,&h(w)=h(b)+1,\\
3,&\text{otherwise.}
\end{cases}
\label{eq:correction}
\end{equation}
Code 3 indexes a sparse stream of the target's high bits. A fused kernel dequantizes
the W4 base, applies the coded or escaped bits, and emits
\begin{equation}
\operatorname{recover}(b,c,m(w))=(h(w)\ll 7)\,|\,m(w)=w.
\label{eq:recovery}
\end{equation}
The equality is at the bit-pattern level, not numerical proximity. Corrections remain
in compact host mappings and pass through reusable pinned staging; only the recovered
expert occupies $C^{16}_{\ell}$ and then executes through the model's native BF16
operator.

\subsection{Unified residency and exact supply}

Every verification request follows one exclusive path:
\begin{equation}
S^{(m)}_{\ell,e}=\begin{cases}
\text{BF16 hit}, & e\in C^{16}_{\ell},\\
\text{exact recover}, & e\notin C^{16}_{\ell},\ m=\mathrm{exact},\\
\text{direct BF16 page-in}, & e\notin C^{16}_{\ell},\ m=\mathrm{ordinary}.
\end{cases}
\label{eq:supply}
\end{equation}
Recovery reuses a resident W4 base or first admits a cold base into $C^4_{\ell}$,
then caches the materialized expert in $C^{16}_{\ell}$. Classic MoE-SubSpec instead
pages the immutable BF16 target after a miss. Requests are deduplicated by
layer/expert identity across the root and proposals, so one recovered expert serves
every routed position in the verification batch.

Representation-aware admission lets a W4 base survive eviction of its materialized
BF16 expert and later regenerate it without another W4 page-in. The controller can
therefore trade several compact bases against fewer ready-to-run exact experts within
the same $B$. Ignoring work shared by both methods, reuse is beneficial when
\begin{equation}
T_{\mathrm{move}}+T_{\mathrm{batch}}>
T_{\mathrm{draft}}+T_{\mathrm{corr}}+T_{\mathrm{recover}}+T_{\mathrm{ctrl}},
\label{eq:benefit}
\end{equation}
where the left side captures avoided expert movement and target work, and the right
side captures proposal, correction, reconstruction, and orchestration. Route overlap,
accepted length, and the W4/BF16 capacity split jointly shape this inequality;
Fig.~\ref{fig:f3} measures each consequence rather than attributing speedup to
compression alone.

\subsection{Shared KV and target authority}

Target prefill establishes clean length $n$ and retains its final input as the first
root. Drafting appends provisional inputs to the same allocation. Verification resets
only logical length, recomputes the root/proposal block with exact experts, and
overwrites that suffix; commit advances through the target-approved prefix, and the
emitted token becomes the next root. No second prefill, duplicate KV cache, or
cross-cache copy is required.

Storage reuse never grants authority to draft state. Executable invariants require an
unchanged physical KV pointer, a byte-identical clean prefix before verification, and
a target-written committed suffix. A proposal is accepted only when it matches the
native target token, with the first mismatch replaced by the target. Together with
Eq.~\eqref{eq:recovery}, native-operator checks, and rejection of undeclared BF16
fallback, these invariants make drafting approximate but verification exact.

\subsection{Round protocol and compositional correctness}

Table~\ref{tab:round-protocol} makes the coupling explicit. The controller does not
launch two models and reconcile their states afterward. It advances one target
sequence through two precision views, changing only the expert representation used
at each phase. Before verification, routed requests from the root and every proposal
are unioned by layer and expert identity. A BF16 hit is consumed immediately; every
miss is resolved by Eq.~\eqref{eq:supply}. The resulting expert view then executes
the unmodified target layer over the entire verification block.

\begin{table}[t]
  \centering
  \caption{One coupled decoding round.}
  \label{tab:round-protocol}
  \footnotesize
  \begin{tabular*}{\columnwidth}{@{\extracolsep{\fill}}cll@{}}
    \toprule
    Step & Active view & State transition \\
    \midrule
    0 & BF16 target & prefill once; establish clean KV \\
    1 & target W4 & draft $\gamma$ tokens in shared KV \\
    2 & controller & rewind logical length; union routes \\
    3 & W4 $+$ correction & materialize requested BF16 misses \\
    4 & BF16 target & verify root and proposals in one pass \\
    5 & controller & commit match or target replacement \\
    \bottomrule
  \end{tabular*}
\end{table}

The protocol composes three local guarantees into an end-to-end one. First,
\emph{representation integrity} follows from Eq.~\eqref{eq:recovery}: a recovered
expert is byte-identical to its checkpoint tensor. Second, \emph{execution integrity}
follows because that tensor enters the native BF16 operator under the target router
and dense path. Third, \emph{sequence integrity} follows because only target-written
KV entries cross the commit boundary. Induction over decoding rounds therefore gives
the same committed token sequence as native target verification under the same
sampling rule. Approximation influences proposal efficiency, never output authority.

This separation also makes the implementation modular. Expert recovery is a supply
operation behind the cache interface; speculative scheduling consumes the resulting
exact view without knowing whether it was a hit, a reconstruction, or an ordinary
page-in. Conversely, the residency controller observes reuse at expert granularity
without interpreting token acceptance. The two modules exchange only requested
expert identities, cache handles, and the clean KV boundary, allowing each to be
tuned while preserving the exactness contract.
